
\begin{lemma}[Direct unary real closure and harmless actual owners]
\label{lem:consolidated-unary-closure}
Suppose $M$ is infinite, with unique minimal infinite invariant
equivalence $E$, and let $F=\Sigma(\Id)$ be its greatest finite-class
equivalence.  Then for every finite real $A$,
\[
                  \operatorname{acl}_M(A)=F[A].
\]
Naming finitely many actual $E$-owner classes does not enlarge
real algebraic closure.  Hence in an owner-named union $N_B$
over any finite $G$-invariant owner set $B$ of a fully normalised
source, real closure is precisely $F[-]$.
\end{lemma}
\begin{proof}
The inclusion $F[A]\subseteq\operatorname{acl}_M(A)$ holds
because $F$ has finite, definable classes.  Conversely, take
$x\notin F[A]$.  For every $a\in A$, the actual invariant
equivalence $\Gamma(\col(x,a))$ has infinite classes: if it
had finite classes, it would lie below the maximal finite-class
invariant equivalence $F$, contradicting $x\notin F[a]$.
Since $E$ is the unique minimal infinite invariant equivalence,
it lies in every such $\Gamma(\col(x,a))$.  Consequently
\[
 E[x]\subseteq
 \bigcap_{a\in A}[a]_{\Gamma(\col(x,a))}.
\]
The left side is infinite.  The contrapositive of the
canonical real-hull theorem \cref{thm:rh-hull} gives
$x\notin\operatorname{acl}_M(A)$.  For $A=\varnothing$,
there are no algebraic real points because $M$ is infinite and
transitive.  Thus $\operatorname{acl}_M(A)=F[A]$.

Since $E$ contains the unique minimal infinite direction,
\cref{lem:acf-harmless-owners} shows that naming any finite
collection $B$ of actual $E$-classes leaves the same real
algebraic closure over $A$.

The owner union $N_B$ is invariant under automorphisms fixing
all named owners, and is homogeneous by joint real/imaginary
exact-sort homogeneity.  If $B$ is $G$-invariant, then
$F[A]\subseteq N_B$ whenever $A\subseteq N_B$; it is
$A$-definable there using the literal finite thin translations.
For $x\in N_B\setminus F[A]$, the preceding harmless-owner
statement makes its orbit under $\Aut(M/A,B)$ infinite.
This group preserves $N_B$ and restricts to a subgroup of
$\Aut(N_B/A)$, so $x$ is not algebraic in $N_B$ over $A$.
Hence $\operatorname{acl}_{N_B}(A)=F[A]$ as claimed.
\end{proof}
